\section{Analyse Muster von Heide Balzert}
\textbf{Heide Balzert} (* 31 Januar 1955 in Essen) ist eine deutsche Informatikerin und Profesorin im Bereich \textbf{Softwaretechnik und Systemanalyse}. In eines ihrer Bücher \textbf{"Lehrbuch der Objektmodellierung"} untersucht sie verschiedene AnalyseMuster und Beispielanwendungen, uns interessiert prinzipiell der Muster nach dem \textbf{Exemplartyp}. Balzert beschreibt den Exemplartyp als
\begin{itemize}
  \item{Einfache Assoziation}
  \item{erstellte Objektverbindungen nicht verändert werden und werden nur gelöscht, wenn das betrefende Exemplar gelöscht wird}
  \item{Name der neuen Klasse enthält oft Begriffe wie Typ, Gruppe und Beschreibung}
\end{itemize}
Dieses Muster untersucht das Problem, dass mehrere Objekten zum gleichen Typ gehören und gemeinsame Eigenschaften besitzen. Laut der \textit{UML-Diagramm in dem Buch} werden verschiedene Buchexemplare dargestellt, die gleiche Buchbeschreibungen haben. Balzert sprach vom Redundante Speicherung von Attributwerten, da die in der Klasse gespeicherten Informationen immer wieder wiederholt werden. Die Abbildung eine zweite Klasse, die die Eigenschaften \textit{Titel, Autor und Verlag} als Attributen enthält während Buchexemplar \cite{Balzert1955} nur Daten für diese Exemplare besitzt, bildet eine Lösung füür das oben genannte Problem. Durch diese Lösung können Redundanzen vermieden werden und somit ist das Modell übersichtlicher.

